% External-artifact case study; numerical values use the frozen case-study macros.
% All numbers are \CS* macros from case_study_macros.tex; do not type numbers here.
% Revised after the 2026-10-19 audit: DEVIATIONS D-006 (Axis I wording), D-007 (reporting-block
% classes and counting), D-008 (scope of the package claim). Previous version: history/paper_d030306/.
\subsection{External case study: a third-party trace through the package}
\label{sec:case-study}

\paragraph{Why this case.}
MoE-Beyond~\cite{gavhane2025moebeyond} is the only study in our audit
(Appendix~A) that publishes both its traces and its trace simulator. It reports
that, for DeepSeek-V2-Lite at batch size one with 10\% of experts resident, a
learned predictor raises the cache hit rate from \CSreportedInfinity{} for the
MoE-Infinity heuristic to \CSreportedBeyond{}. Setting that comparison beside
other studies requires fields its public text does not state: event replay,
cache scope and capacity, hit counting, and workload construction. This section
applies the released package, with documented analyst steps, to the public
artifact (commit \texttt{\CSthirdPartyCommit}) to answer those fields. It does
not re-evaluate the reported result. The analysis was pre-registered before any
result was computed; every deviation, including corrections made after an
independent audit, is listed with the artifact.

\paragraph{Execution contract.}
Running the published artifact at the reported configuration reproduced the
reported value (\CSreproInfinity{} vs \CSreportedInfinity{}). Running the
published artifact at the reported configuration reproduced the reported value
(\CSreproBeyond{} vs \CSreportedBeyond{}). Both are means over \CSseeds{} seeds
of the random cache pre-fill (tolerance $\pm$\CSreproTolerance{} percentage
points). On the in-event eviction probe, the published simulator recorded
\CSprobeAtwoMisses{} misses, the count expected under per-access retention; its
executed replay contract is classified as per-access. On the reference trace of
Section~2.5 the published simulator recorded \CSprobeAoneHits{} hits and
\CSprobeAoneMisses{} misses, matching the event-atomic reference counts
(12 / 18). The simulator's executed contract is per-access (probe A2); its
reported counts are per-access counts, and setting them beside an event-atomic
study requires declaring this (checklist A1--A2). The traces record one prefill
forward pass per prompt, so we declare two event definitions: per-token events,
matching the simulator's replay unit, and one whole-prefill event per layer,
matching the executed forward pass. At the reported operating point (batch size
one, 10\% residency, global scope, \CStheirCapacity{} blocks) under either event
definition, every event's expert set fits within capacity. This is a capacity
observation; it does not by itself establish equal counts under per-access and
event-atomic replay, because a per-access miss can evict a resident expert
needed later in the same event. A two-event trace in the package's regression
tests, in which both events fit, is counted differently by the two replay
contracts. At our per-layer frame with whole-prefill events, a fraction
\CSoverCapEtwoPerLayerTen{} of events exceed capacity; in that setting the
replay contract can affect the counts and should be declared (checklist A1--A2).

\paragraph{Operating regime.}
Under our reporting frame, the per-token event stream of this setting has
$\bar r = \CSrbarEoneRhoForty$ at $\rho = 40\%$ (per-layer scope); the nearest
measured alignment point in Table~11 is the \CSnearestEone{}
($\bar r = \CSnearestEoneRbar$). Under our reporting frame, the whole-prefill
event stream of this setting has $\bar r = \CSrbarEtwoRhoForty$ at $\rho = 40\%$
(per-layer scope); the nearest measured alignment point in Table~11 is the
\CSnearestEtwo{} ($\bar r = \CSnearestEtwoRbar$).

\paragraph{Workload construction and a package-contract analysis.}
The \CStestPrompts{} public test prompts have \CSheadsDistinct{} distinct
70-character heads, \CSfinalsDistinct{} distinct final lines and a 95th-percentile
word 5-gram Jaccard of \CSpromptJaccardPninetyfive{}; the collection scripts
apply no chat template and record no decode step. On the public test prompts, the \S5.8 detector reports
flag~=~\CSdetectorFlag{}; this is a detector-level screening result, not a
replication or refutation of any workload effect in that study. The triggering
positional agreement sits at position~0, a shared beginning-of-sequence token
(\CSonlyBosShare{} of pairs agree only there), and at the question opener, in
sequences of \CSlenMin{}--\CSlenMax{} tokens; the threshold was validated on
512-token sequences. Excluding position~0 post hoc lowers the 95th percentile
from \CSposAgreePninetyfive{} to \CSposAgreeNoBosPninetyfive{}. This is a
limitation of our own detector: positional agreement is length-sensitive and
needs a length-aware threshold.
Separately, and under the package's own contract rather than the published one
(event-atomic, prefetch-free, cold start; per-token events, per-layer scope,
$\rho = 40\%$), the offline optimum leaves a recoverable gap of \CSgapRhoForty{},
of which the future-victim share is \CSfvShareRhoForty{}. These are
prefetch-free residency results under our replay contract and are not comparable
with the prefetch-inclusive hit rates reported by MoE-Beyond. The same analysis
bounds the decomposition of Section~7.2: with whole-prefill events at
$\rho = 40\%$ ($\bar r = \CSrbarEtwoRhoForty$), forced admission performs worse
than the best causal policy and the two shares leave $[0,1]$ (\CSadmShareBreak{}
and \CSfvShareBreak{}); the decomposition presupposes $m_B \le m_F \le m_c$,
which no longer holds once most events exceed their quota.

\paragraph{What the workflow added.}
Each reporting-block field is classified for the published configuration under
one rule for both sources: answered, established as not applicable, not
available, or declarable only by the original authors. From the paper text
alone, \CSfieldsPaperOnly{} of the \CSfieldsTotal{} reporting-block fields are
answered (\CSfieldsPaperOnlyNotApplicable{} more are established as not
applicable); with the public artifact, the package and the documented analyst
steps, \CSfieldsWithArtifacts{} are answered
(\CSfieldsWithArtifactsNotApplicable{} not applicable). Of the
\CSfieldsAnsweredOnlyWithArtifacts{} fields answered only in the second mode,
\CSaddedByExecution{} come from executing the package's diagnostics or the
published simulator, \CSaddedByCodeInspection{} from reading the published code
and \CSaddedByDataInspection{} from inspecting the published data; the per-field
provenance is in the artifact. The comparison has no artifact-only control, so
the difference is a property of the combined workflow, not of the package alone.
What the package contributed is concrete: conformance probes in the third-party
input format, regime and workload diagnostics on the normalized trace, and a
reporting block whose every value carries its class, provenance and input hash.
